%%%%%%%%%%%%%%%%%%%%%%%%%%%%%%%%%%%%%%%%%%%%%%%%%%%%%%%%%%%%%%%%%%%%%%%%%%%
%%  appendixC.tex  –  Дополнительные экспериментальные результаты
%%%%%%%%%%%%%%%%%%%%%%%%%%%%%%%%%%%%%%%%%%%%%%%%%%%%%%%%%%%%%%%%%%%%%%%%%%%

\chapter{Дополнительные экспериментальные результаты}
\label{app:extra_results}

В данном приложении представлены дополнительные экспериментальные результаты, дополняющие анализ Главы~\ref{ch:experiments} (Экспериментальные результаты и валидация). Раздел~\ref{app:per_attack} предоставляет покатегорийные разбивки производительности по типам атак. Раздел~\ref{app:ablation_detail} документирует детали покомпонентной абляции и исследования чувствительности к гиперпараметрам. Раздел~\ref{app:cross_cloud} представляет матрицы кросс-облачного и кросс-доменного переноса. Раздел~\ref{app:privacy_utility} анализирует компромисс приватность--качество детально. Раздел~\ref{app:comp_perf} представляет бенчмарки вычислительной производительности и характеристики масштабирования.


%%=====================================================================
\section{Производительность по категориям атак}
\label{app:per_attack}
%%=====================================================================

\subsection{Детальные результаты CIC-IoT-2023}
\label{app:cicids_iot_detail}

Таблица~\ref{tab:app_iot_peratk} представляет покатегорийные прецизионность, полноту, F1 и AUC-ROC для CT-TGNN на CIC-IoT-2023. Все семь основных категорий превышают 95\% полноты, при этом DDoS достигает наивысшей полноты 99,1\% благодаря характерной объёмной сигнатуре, а атаки спуфинга демонстрируют наименьшую полноту 95,2\%, поскольку протокольная имитация может близко воспроизводить легитимные паттерны аутентификации.

\begin{table}[htbp]
\centering
\caption{Покатегорийные результаты CT-TGNN на CIC-IoT-2023.}
\label{tab:app_iot_peratk}
\begin{tabular}{lcccc}
\toprule
Категория атак & Прецизионность & Полнота & F1 & AUC-ROC \\
\midrule
DDoS             & 98,7\% & 99,1\% & 98,9\% & 0,998 \\
Разведка         & 96,2\% & 96,8\% & 96,5\% & 0,991 \\
Веб-атаки        & 97,1\% & 97,3\% & 97,2\% & 0,993 \\
Перебор           & 98,4\% & 98,6\% & 98,5\% & 0,997 \\
Спуфинг          & 94,8\% & 95,2\% & 95,0\% & 0,987 \\
ВПО              & 98,6\% & 98,9\% & 98,7\% & 0,998 \\
Варианты Mirai   & 97,1\% & 97,4\% & 97,2\% & 0,994 \\
\midrule
Легитимный       & 98,9\% & 98,4\% & 98,7\% & 0,997 \\
\bottomrule
\end{tabular}
\end{table}


\subsection{Детальные результаты CSE-CICIDS2018}
\label{app:cicids2018_detail}

Таблица~\ref{tab:app_cicids_peratk} представляет покатегорийную производительность единой системы на CSE-CICIDS2018. Единая система достигает наибольшего улучшения над базовыми линиями на инфильтрации (12,7 п.\,п.\ выше лучшей базовой линии), обнаружении ботнетов (9,3 п.\,п.) и эксплуатации Heartbleed (8,9 п.\,п.) --- категориях, включающих многоэтапные, темпорально протяжённые последовательности атак, где формулировка непрерывного времени обеспечивает наибольшее преимущество.

\begin{table}[htbp]
\centering
\caption{Покатегорийные результаты единой системы на CSE-CICIDS2018.}
\label{tab:app_cicids_peratk}
\begin{tabular}{lccccc}
\toprule
Категория атак & Прециз. & Полнота & F1 & AUC & $\Delta$ vs.\ Базов. \\
\midrule
Перебор (SSH)        & 95,3\% & 96,1\% & 95,7\% & 0,991 & +3,8\% \\
Перебор (FTP)        & 96,7\% & 97,3\% & 97,0\% & 0,994 & +4,2\% \\
DoS Hulk             & 97,8\% & 98,2\% & 98,0\% & 0,997 & +2,1\% \\
DoS GoldenEye        & 96,4\% & 96,9\% & 96,6\% & 0,993 & +3,4\% \\
DoS Slowloris        & 95,1\% & 95,7\% & 95,4\% & 0,989 & +5,1\% \\
Веб-атака (XSS)      & 93,8\% & 94,3\% & 94,1\% & 0,984 & +6,7\% \\
Веб-атака (SQL Inj.) & 94,2\% & 94,8\% & 94,5\% & 0,986 & +5,9\% \\
Инфильтрация         & 91,7\% & 92,4\% & 92,0\% & 0,978 & +12,7\% \\
Ботнет (Ares)        & 93,1\% & 93,8\% & 93,4\% & 0,982 & +9,3\% \\
Heartbleed           & 92,3\% & 93,1\% & 92,7\% & 0,979 & +8,9\% \\
Сканирование портов  & 97,6\% & 98,1\% & 97,8\% & 0,996 & +2,3\% \\
DDoS LOIC            & 98,1\% & 98,4\% & 98,2\% & 0,998 & +1,7\% \\
\midrule
Легитимный           & 97,4\% & 96,8\% & 97,1\% & 0,995 & +2,4\% \\
\bottomrule
\end{tabular}
\end{table}


\subsection{Детальные результаты UNSW-NB15}
\label{app:unsw_detail}

Таблица~\ref{tab:app_unsw_peratk} представляет покатегорийные результаты на UNSW-NB15. Единая система достигает наибольшего улучшения на атаках разведки (8,7 п.\,п.), где многомасштабное темпоральное моделирование захватывает медленные и скрытные паттерны сканирования, которые базовые линии дискретного времени систематически пропускают.

\begin{table}[htbp]
\centering
\caption{Покатегорийные результаты единой системы на UNSW-NB15.}
\label{tab:app_unsw_peratk}
\begin{tabular}{lccccc}
\toprule
Категория атак & Прециз. & Полнота & F1 & AUC & $\Delta$ vs.\ Базов. \\
\midrule
Generic          & 96,8\% & 97,2\% & 97,0\% & 0,995 & +3,2\% \\
Exploits         & 94,7\% & 96,3\% & 95,5\% & 0,989 & +5,4\% \\
Fuzzers          & 93,2\% & 94,1\% & 93,6\% & 0,984 & +4,8\% \\
DoS              & 96,4\% & 97,1\% & 96,7\% & 0,993 & +3,7\% \\
Разведка         & 92,8\% & 97,9\% & 95,3\% & 0,988 & +8,7\% \\
Analysis         & 97,1\% & 98,4\% & 97,7\% & 0,996 & +4,1\% \\
Backdoors        & 91,9\% & 93,7\% & 92,8\% & 0,981 & +6,3\% \\
Shellcode        & 94,3\% & 95,8\% & 95,0\% & 0,987 & +5,1\% \\
Worms            & 89,7\% & 91,4\% & 90,5\% & 0,973 & +7,2\% \\
\midrule
Нормальный       & 97,2\% & 96,4\% & 96,8\% & 0,994 & +3,4\% \\
\bottomrule
\end{tabular}
\end{table}


\subsection{Разбивка атак микросервисов TripleE-TGNN}
\label{app:triplee_detail}

Таблица~\ref{tab:app_triplee_peratk} представляет покатегорийную точность обнаружения TripleE-TGNN на наборе данных Container Security, разбитую по гранулярности, на которой каждая атака наиболее эффективно обнаруживается.

\begin{table}[htbp]
\centering
\caption{Покатегорийная точность обнаружения TripleE-TGNN и доминирующая гранулярность.}
\label{tab:app_triplee_peratk}
\begin{tabular}{lccc}
\toprule
Тип атаки & Точность & Домин. гранулярность & Одногран. точн. \\
\midrule
Побег из контейнера  & 98,7\% & Узел + Сервис & 91,2\% \\
Цепь поставок        & 97,3\% & Сервис + Трассировка & 89,8\% \\
Исчерпание ресурсов  & 96,9\% & Сервис + Узел & 90,3\% \\
Повышение привилегий & 97,8\% & Узел & 93,1\% \\
Злоупотребление API  & 95,4\% & Сервис & 91,7\% \\
Латер. перемещение   & 96,1\% & Трассировка + Узел & 88,4\% \\
Эксфильтрация данных & 94,8\% & Трассировка & 87,6\% \\
Криптоджекинг        & 97,2\% & Узел & 92,9\% \\
\midrule
Общий результат      & 96,8\% & Многоуровневый & 88,5\% \\
\bottomrule
\end{tabular}
\end{table}


%%=====================================================================
\section{Детали абляционного исследования}
\label{app:ablation_detail}
%%=====================================================================

\subsection{Покомпонентная абляция с доверительными интервалами}
\label{app:component_ablation}

Таблица~\ref{tab:app_ablation_ci} расширяет результаты абляции из Таблицы~\ref{tab:ablation} основного текста 95\% доверительными интервалами, вычисленными по пяти независимым запускам с различными случайными затравками. Все деградации статистически значимы после коррекции Бонферрони на уровне $\alpha=0.05$.

\begin{table}[htbp]
\centering
\caption{Расширенные результаты абляции с 95\% доверительными интервалами (пять затравок).}
\label{tab:app_ablation_ci}
\begin{tabular}{lccc}
\toprule
Удалённый компонент & Полная система & Без компонента & $d$ Коэна \\
\midrule
НВ ОДУ                   & $98.3{\pm}0.2$\% & $91.5{\pm}0.4$\% & 4,82 \\
Многомасшт. моделиров.    & $98.3{\pm}0.2$\% & $94.4{\pm}0.3$\% & 3,14 \\
Темпор. адапт. БН         & $98.3{\pm}0.2$\% & $93.6{\pm}0.5$\% & 3,67 \\
Межгранул. внимание       & $96.8{\pm}0.3$\% & $94.7{\pm}0.3$\% & 2,41 \\
Усеч. среднее $\to$ FedAvg & $89.7{\pm}0.4$\% & $68.3{\pm}1.2$\% & 7,13 \\
Дифф. приватность          & $87.1{\pm}0.3$\% & $90.2{\pm}0.3$\% & $-$3,28 \\
Выравнивание ОТ            & $92.7{\pm}0.4$\% & $76.9{\pm}0.8$\% & 5,92 \\
Ветвь LLM (F1)            & $87.6{\pm}0.5$\% & $34.2{\pm}1.1$\% & 8,47 \\
Прогресс. дистилляция      & $91.4{\pm}0.3$\% & $81.7{\pm}0.6$\% & 4,71 \\
Вариац. внимание (ECE)     & $0.078{\pm}0.004$ & $0.192{\pm}0.008$ & 5,38 \\
Теор.-игровое равновесие   & $94.7{\pm}0.3$\% & $87.6{\pm}0.4$\% & 4,18 \\
\bottomrule
\end{tabular}
\end{table}


\subsection{Анализ чувствительности к гиперпараметрам}
\label{app:sensitivity}

В данном разделе представлены анализы чувствительности для наиболее важных гиперпараметров каждого метода, при варьировании одного параметра при фиксации всех остальных на значениях по умолчанию.

\subsubsection{CT-TGNN: Допуск решателя ОДУ}

Таблица~\ref{tab:app_sens_tol} представляет влияние допуска решателя ОДУ на точность, задержку вывода и число вычислений функции (NFE) на прямой проход. Более жёсткие допуски дают маргинальное улучшение точности, но существенно увеличивают вычислительную стоимость; допуск по умолчанию $10^{-5}$ обеспечивает лучший компромисс точность--эффективность.

\begin{table}[htbp]
\centering
\caption{Чувствительность CT-TGNN к относительному допуску решателя ОДУ.}
\label{tab:app_sens_tol}
\begin{tabular}{lccc}
\toprule
Отн. допуск & Точность & Задержка (мс) & Ср. NFE \\
\midrule
$10^{-3}$ & 97,6\% & 28 & 14 \\
$10^{-4}$ & 98,1\% & 36 & 21 \\
$10^{-5}$ (по умолч.) & 98,3\% & 47 & 31 \\
$10^{-6}$ & 98,4\% & 72 & 48 \\
$10^{-7}$ & 98,4\% & 118 & 73 \\
\bottomrule
\end{tabular}
\end{table}


\subsubsection{FedLLM-API: Бюджет приватности и доля усечения}

Таблица~\ref{tab:app_sens_dp} представляет влияние варьирования бюджета приватности на раунд~$\epsilon$ на точность, общий расход приватности за 100~раундов и успех атаки вывода принадлежности. Значение по умолчанию $\epsilon=0.85$ балансирует строгую приватность с допустимой деградацией точности.

\begin{table}[htbp]
\centering
\caption{Чувствительность FedLLM-API к бюджету дифференциальной приватности $\epsilon$.}
\label{tab:app_sens_dp}
\begin{tabular}{lcccc}
\toprule
$\epsilon$ & Точность & Общий $\epsilon_{100}$ & Атака ВП & $\Delta$ Точн. vs.\ Без-ДП \\
\midrule
0,5  & 84,2\% & 5,0  & 51,3\% & $-$7,0\% \\
0,85 (по умолч.) & 87,1\% & 8,5  & 52,7\% & $-$3,1\% \\
1,0  & 88,4\% & 10,0 & 54,1\% & $-$1,8\% \\
1,5  & 89,7\% & 15,0 & 57,8\% & $-$0,5\% \\
2,0  & 91,8\% & 20,0 & 62,4\% & +1,6\% \\
$\infty$ (без ДП) & 90,2\% & $\infty$ & 71,3\% & 0\% \\
\bottomrule
\end{tabular}
\end{table}

Таблица~\ref{tab:app_sens_trim} представляет точность при 30\% византийских противников при варьировании доли усечения. Доли усечения ниже 20\% не удаляют все византийские обновления, а доли выше 30\% отбрасывают слишком много честных обновлений.

\begin{table}[htbp]
\centering
\caption{Чувствительность FedLLM-API к доле усечения при 30\% византийских противников.}
\label{tab:app_sens_trim}
\begin{tabular}{lcc}
\toprule
Доля усечения & Точность (30\% визант.) & Раунды сходим. \\
\midrule
10\% & 72,4\% & 142 \\
15\% & 81,3\% & 108 \\
20\% (по умолч.) & 89,7\% & 78 \\
25\% & 88,9\% & 82 \\
30\% & 86,1\% & 91 \\
35\% & 82,7\% & 104 \\
\bottomrule
\end{tabular}
\end{table}


\subsubsection{MambaShield: Температура дистилляции и размер ансамбля}

Таблица~\ref{tab:app_sens_mamba} представляет совместное влияние температуры дистилляции и размера ансамбля учителей на робастную точность при 25\% отравления. Температура~3.0 с пятью учителями даёт лучшую производительность; более низкие температуры порождают более резкие распределения, усиливающие шум меток, а бо\'{л}ьшие ансамбли дают убывающую отдачу относительно их вычислительной стоимости.

\begin{table}[htbp]
\centering
\caption{Чувствительность MambaShield к температуре дистилляции и размеру ансамбля (25\% отравления).}
\label{tab:app_sens_mamba}
\begin{tabular}{lcccc}
\toprule
Температура & 3 учит. & 5 учит. & 7 учит. & 10 учит. \\
\midrule
$T=1.0$ & 85,3\% & 86,7\% & 87,1\% & 87,3\% \\
$T=2.0$ & 88,4\% & 89,8\% & 90,1\% & 90,2\% \\
$T=3.0$ (по умолч.) & 89,7\% & 91,4\% & 91,6\% & 91,7\% \\
$T=5.0$ & 88,1\% & 90,2\% & 90,5\% & 90,6\% \\
$T=10.0$ & 84,6\% & 87,3\% & 87,7\% & 87,8\% \\
\bottomrule
\end{tabular}
\end{table}


\subsubsection{Стохастический трансформер: Выборки Монте-Карло}

Таблица~\ref{tab:app_sens_mc} представляет влияние числа прямых проходов Монте-Карло на ECE, корреляцию эпистемической неопределённости с ошибкой предсказания и задержку вывода на батч. Ранняя остановка по сходимости дисперсии снижает эффективное число проходов до 18 в среднем.

\begin{table}[htbp]
\centering
\caption{Чувствительность стохастического трансформера к числу прямых проходов МК.}
\label{tab:app_sens_mc}
\begin{tabular}{lcccc}
\toprule
Проходов МК & ECE & Корр. неопред. & Задержка (мс/батч) & Эфф. проходов \\
\midrule
5   & 0,124 & 0,43 & 78  & 5 \\
10  & 0,103 & 0,52 & 126 & 10 \\
20  & 0,089 & 0,61 & 204 & 17 \\
50 (по умолч.) & 0,078 & 0,67 & 340 & 18 \\
100 & 0,076 & 0,68 & 638 & 19 \\
\bottomrule
\end{tabular}
\end{table}


%%=====================================================================
\section{Кросс-облачный и кросс-доменный перенос}
\label{app:cross_cloud}
%%=====================================================================

\subsection{Матрица кросс-датасетного переноса}
\label{app:transfer_matrix}

Таблица~\ref{tab:app_transfer} представляет точность CT-TGNN при обучении на одном наборе данных (строки) и тестировании на другом (столбцы) без дообучения. Диагональные записи соответствуют производительности внутри распределения. Стабильно малая внедиагональная деградация, в среднем 3,8\% по всем парам, подтверждает свойства обобщения формулировки темпоральных графов непрерывного времени.

\begin{table}[htbp]
\centering
\caption{Точность кросс-датасетного переноса CT-TGNN (\%). Строки: обучающий набор; столбцы: тестовый набор.}
\label{tab:app_transfer}
\begin{tabular}{lccc}
\toprule
Обуч. $\backslash$ Тест & CIC-IoT & CICIDS18 & UNSW-NB15 \\
\midrule
CIC-IoT-2023   & 98,3 & 92,7 & 91,4 \\
CSE-CICIDS2018 & 91,8 & 96,7 & 89,3 \\
UNSW-NB15      & 90,6 & 88,9 & 97,1 \\
\bottomrule
\end{tabular}
\end{table}


\subsection{Кросс-доменный перенос на речь и здравоохранение}
\label{app:cross_domain_detail}

Таблица~\ref{tab:app_cross_domain} представляет детальные метрики для экспериментов кросс-доменного переноса, описанных в Разделе~\ref{sec:transfer} основного текста, включая влияние предобучения на сетевой безопасности.

\begin{table}[htbp]
\centering
\caption{Результаты кросс-доменного переноса с и без предобучения на сетевой безопасности.}
\label{tab:app_cross_domain}
\begin{tabular}{llcccc}
\toprule
Домен & Метод & Точность & F1 & ECE & Предобуч.? \\
\midrule
Речь (LibriSpeech) & CT-TGNN & 90,4\% & 90,8\% & 0,091 & Нет \\
Речь (LibriSpeech) & CT-TGNN & 94,2\% & 94,2\% & 0,073 & Да \\
Речь (LibriSpeech) & LSTM базовая & 86,9\% & 86,7\% & 0,142 & Нет \\
\midrule
Здравоохранение (MIMIC-III) & Стох.\ трансф. & 86,1\% & 85,7\% & 0,098 & Нет \\
Здравоохранение (MIMIC-III) & Стох.\ трансф. & 89,7\% & 89,3\% & 0,082 & Да \\
Здравоохранение (eICU) & Стох.\ трансф. & 84,8\% & 84,2\% & 0,104 & Нет \\
Здравоохранение (eICU) & Стох.\ трансф. & 88,4\% & 87,9\% & 0,087 & Да \\
\bottomrule
\end{tabular}
\end{table}


%%=====================================================================
\section{Анализ компромисса приватность--качество}
\label{app:privacy_utility}
%%=====================================================================

\subsection{Влияние дифференциальной приватности по наборам данных}
\label{app:dp_impact}

Таблица~\ref{tab:app_dp_datasets} представляет деградацию точности, вызванную дифференциальной приватностью при $\epsilon=0.85$, по всем трём основным наборам данных. Деградация стабильна по наборам данных, составляя от 2,8\% до 3,4\%, подтверждая, что компромисс приватность--качество не варьируется существенно с распределением данных.

\begin{table}[htbp]
\centering
\caption{Влияние дифференциальной приватности на точность по наборам данных ($\epsilon=0.85$, $\delta=10^{-5}$).}
\label{tab:app_dp_datasets}
\begin{tabular}{lccc}
\toprule
Набор данных & Без ДП & С ДП & Деградация \\
\midrule
CIC-IoT-2023   & 94,7\% & 91,6\% & $-$3,1\% \\
CSE-CICIDS2018 & 93,9\% & 91,1\% & $-$2,8\% \\
UNSW-NB15      & 93,8\% & 90,4\% & $-$3,4\% \\
\bottomrule
\end{tabular}
\end{table}


\subsection{Учёт приватности по раундам коммуникации}
\label{app:privacy_accounting}

Таблица~\ref{tab:app_privacy_rounds} представляет кумулятивный расход приватности как функцию федеративных раундов коммуникации, вычисленный через моменты-бухгалтер с усилением пуассоновской подвыборкой. Общий бюджет $\epsilon_{\mathrm{total}}=8.5$ после 100~раундов остаётся в диапазоне, обычно считающемся приемлемым для чувствительных приложений.

\begin{table}[htbp]
\centering
\caption{Кумулятивный расход приватности по федеративным раундам ($\delta=10^{-5}$).}
\label{tab:app_privacy_rounds}
\begin{tabular}{lcccccc}
\toprule
Раунды & 10 & 25 & 50 & 75 & 100 & 150 \\
\midrule
$\epsilon_{\mathrm{total}}$ & 0,85 & 2,13 & 4,25 & 6,38 & 8,50 & 12,75 \\
Точность & 78,4\% & 84,1\% & 87,8\% & 88,6\% & 87,1\% & 86,3\% \\
\bottomrule
\end{tabular}
\end{table}


%%=====================================================================
\section{Метрики вычислительной производительности}
\label{app:comp_perf}
%%=====================================================================

\subsection{Пропускная способность и задержка вывода}
\label{app:throughput}

Таблица~\ref{tab:app_throughput} сравнивает пропускную способность вывода, задержку на образец, потребление памяти GPU и энергоэффективность по всем семи методам и базовым линиям. Измерения проведены на одиночном GPU NVIDIA Tesla P100 с размером батча~32 для задержки и размером батча~512 для пропускной способности, представляя медиану 100~запусков с 95\% бутстрап-доверительными интервалами.

\begin{table}[htbp]
\centering
\caption{Сравнение вычислительной производительности по методам (P100 GPU).}
\label{tab:app_throughput}
\begin{tabular}{lcccc}
\toprule
Метод & Пропуск. (тыс./с) & Задержка (мс) & Память (ГБ) & Вт/1 тыс. предск. \\
\midrule
CT-TGNN            & 8{,}700 & 47 & 3,8 & 34 \\
TripleE-TGNN       & 4{,}200 & 73 & 4,7 & 51 \\
FedLLM-API         & 5{,}100 & 62 & 5,2 & 47 \\
Прямо совм. обнаруж. & 6{,}800 & 38 & 2,1 & 28 \\
MambaShield        & 12{,}400 & 18 & 2,3 & 19 \\
Стох.\ трансформер  & 2{,}300 & 340 & 8,2 & 89 \\
Теор.-игровой      & 7{,}100 & 12 & 0,2 & 8 \\
\midrule
LSTM базовая       & 4{,}300 & 38 & 1,8 & 42 \\
Стд. трансформер   & 6{,}200 & 55 & 4,1 & 125 \\
Random Forest      & 14{,}200 & 12 & 0,8 & 14 \\
XGBoost            & 18{,}700 & 8 & 0,6 & 11 \\
\bottomrule
\end{tabular}
\end{table}


\subsection{Масштабирование с размером входа}
\label{app:scaling}

Таблица~\ref{tab:app_scaling} представляет зависимость задержки вывода от длины входной последовательности для CT-TGNN, MambaShield и стандартного трансформера. CT-TGNN масштабируется линейно через фиксированное число шагов решателя ОДУ на постоянную времени. MambaShield масштабируется линейно через механизм селективного сканирования. Стандартный трансформер масштабируется квадратично через самовнимание, становясь непрактичным для последовательностей длиннее приблизительно 4096~токенов на 16\,ГБ GPU.

\begin{table}[htbp]
\centering
\caption{Задержка вывода (мс) как функция длины последовательности.}
\label{tab:app_scaling}
\begin{tabular}{lccccc}
\toprule
Длина последоват. & 256 & 512 & 1024 & 2048 & 4096 \\
\midrule
CT-TGNN          & 23 & 41 & 74 & 142 & 278 \\
MambaShield      & 8  & 14 & 26 & 49  & 94  \\
Стд. трансформер  & 18 & 52 & 178 & 684 & OOM \\
\bottomrule
\end{tabular}
\end{table}


\subsection{Время обучения и сходимость}
\label{app:training_time}

Таблица~\ref{tab:app_training} представляет общее время обучения, число эпох до сходимости и время на эпоху для каждого метода на CIC-IoT-2023 с использованием четырёх V100 GPU с параллелизмом по данным.

\begin{table}[htbp]
\centering
\caption{Сравнение времени обучения (четыре V100 GPU, CIC-IoT-2023).}
\label{tab:app_training}
\begin{tabular}{lccc}
\toprule
Метод & Общее время (ч) & Эпохи & Время/Эпоху (мин) \\
\midrule
CT-TGNN            & 18,2 & 142 & 7,7 \\
TripleE-TGNN       & 22,7 & 118 & 11,5 \\
FedLLM-API         & 31,4 & 100 (раундов) & 18,8 \\
Прямо совм. обнаруж. & 8,3  & 87  & 5,7 \\
MambaShield        & 14,1 & 100 & 8,5 \\
Стох.\ трансформер  & 24,6 & 96  & 15,4 \\
Теор.-игровой      & 6,8  & 10{,}000 (раундов) & 0,04 \\
\bottomrule
\end{tabular}
\end{table}


\subsection{Сравнение размеров моделей}
\label{app:model_size}

Таблица~\ref{tab:app_params} представляет общее число обучаемых параметров, размер сжатой модели после INT8 квантования и точность, сохранённую после квантования.

\begin{table}[htbp]
\centering
\caption{Размер модели и влияние квантования.}
\label{tab:app_params}
\begin{tabular}{lcccc}
\toprule
Метод & Параметры (M) & FP32 (МБ) & INT8 (МБ) & INT8 Точн. \\
\midrule
CT-TGNN            & 12,3 & 47  & 13 & 97,9\% \\
TripleE-TGNN       & 18,7 & 71  & 19 & 96,2\% \\
FedLLM-API         & 24,1 & 92  & 25 & 86,4\% \\
Прямо совм. обнаруж. & 6,8  & 26  & 7  & 95,8\% \\
MambaShield        & 8,4  & 32  & 9  & 90,9\% \\
Стох.\ трансформер  & 31,2 & 119 & 31 & 93,4\% \\
\midrule
Стд. трансформер   & 47,8 & 182 & 47 & 90,8\% \\
\bottomrule
\end{tabular}
\end{table}

